\chapter{Endpoint Threat Auditing}
\label{ch:4.4}
\chaptertopic{Topic 4.4}

\chapterquote{It is a capital mistake to theorize before one has all the evidence.}{Arthur Conan Doyle, \emph{A Study in Scarlet} (1887)}

\section{Motivation}

Chapter~\ref{ch:4.3} hardened the device, watched it continuously, and patched its known flaws.
This closing chapter of Unit~4 asks the question those controls cannot fully answer on their own:
given everything that happened on a device, was it actually compromised? Three complementary
disciplines answer this together -- live behavioral monitoring on the device itself, correlating
historical log entries across time, and combining scattered forensic clues into one confident
verdict. None of the three is sufficient alone; each catches what the others miss.

\begin{apalignment}{Unit 4: Securing Devices}{Topic 4.4 -- Detecting Attacks on Devices}
\begin{apobjectives}
\apobjective{4.4.A}{Detect Attacks}{Account for how an attack on a device is noticed at all: what
the operating system and its security tooling record, why that record has to leave the machine
before an attacker can reach it, and what makes a sequence of events evidence when no single one of
them is.}
\apobjective{4.4.B}{Detect Attacks}{Choose the detective controls that suit a given device and a
given threat -- audit policy and the events it turns on, off-host forwarding, correlation logic,
tamper-evident retention -- and be able to say which attack each one would and would not catch.}
\apobjective{4.4.C}{Detect Attacks}{Judge a detection method by what it does in practice: the
coverage it provides, the volume it generates, whether the people receiving its output can act on
it, and the failure it will still miss -- recognizing that a rule nobody can triage lowers detection
capability rather than raising it.}
\apobjective{4.4.D}{Detect Attacks}{Work through real authentication logs and identify the traces
of a credential attack, separating exhaustive guessing against one account from low-and-slow
guessing across many, and stating which distinctions the log can support and which it cannot.}
\end{apobjectives}
\end{apalignment}

\section{Endpoint Detection and Response}

A firewall guards the network boundary, but once malicious code is already running inside a
machine, a boundary control cannot see it at all.

\begin{definition}{Endpoint Detection and Response (EDR)}
Security software that runs as an agent on a device, continuously collecting telemetry -- process,
file system, registry, network, and user-account events -- and using both signature matching and
behavioral modeling to detect and respond to threats on the endpoint itself.
\end{definition}

The telemetry an EDR agent records is dense. It includes process events (every program starting,
stopping, or spawning a child), file system events (creation, modification, deletion), and registry
changes on Windows (a common mechanism for \textbf{persistence}, letting malware restart after a
reboot). It also includes network connections (which process opened a socket, to where, and how
much data moved) and user-account events (logins, privilege escalation). \textbf{Signature matching}
compares this stream against known-bad indicators, fast and precise for cataloged threats but blind
to anything novel, exactly as in Chapter~\ref{ch:4.3}. \textbf{Behavioral detection} instead builds a
model of what normal looks like on a specific host and flags deviations from it. That catches not
just novel malware but \textbf{living-off-the-land} techniques, where an attacker misuses entirely
legitimate tools like PowerShell or WMI to carry out malicious actions. No file-based signature would
ever flag those actions, because the tool itself is not malware.

\begin{figure}[h]
\centering
\begin{tikzpicture}[
  src/.style={draw=kickstartermaroon, thick, rounded corners, minimum width=2.5cm, minimum height=0.8cm, align=center, font=\scriptsize},
]
\node[src, fill=kickstartercream] (proc) at (0,1.5) {Process events};
\node[src, fill=kickstartercream] (file) at (0,0.5) {File system events};
\node[src, fill=kickstartercream] (reg) at (0,-0.5) {Registry changes};
\node[src, fill=kickstartercream] (net) at (0,-1.5) {Network connections};
\node[src, fill=kickstartermaroon!20] (model) at (4.2,0) {Behavioral model\\(what's normal\\on \emph{this} host)};
\node[src, fill=kickstartermaroon!30] (verdict) at (8.2,0) {Combined score\\$\to$ alert / isolate /\\kill / quarantine};
\foreach \s in {proc, file, reg, net}
  \draw[-{Latex[length=2mm]}, kickstartermaroon] (\s) -- (model);
\draw[-{Latex[length=2mm]}, kickstartermaroon, thick] (model) -- (verdict);
\node[font=\scriptsize\sffamily, gray, align=center] at (4.2,1.6) {no single stream is damning alone --\\the combination is what scores as high-risk};
\end{tikzpicture}
\caption{EDR telemetry sources feed one behavioral model per host; the response action follows from
the combined score, not from any single stream in isolation.}
\label{fig:4.4-edr-telemetry}
\end{figure}

\begin{example}[A living-off-the-land chain]
An employee opens an email attachment whose macro uses PowerShell -- a legitimate Windows tool --
to download a second-stage payload. Traditional antivirus misses this entirely, since PowerShell
itself is not malicious. An EDR agent instead records the full chain: an unusual parent-child
relationship (the Office application spawned PowerShell) and an unusual outbound network connection
from a process that normally never touches the network. It also records a new executable written to
a temp directory. No single event in that chain is damning on its own, but the behavioral model scores the
\emph{combination} as high-risk and can isolate the laptop automatically before the payload ever
reaches another system.
\end{example}

Detection is only the first step. EDR is distinguished by its \textbf{response} capability once
confidence is high enough: it can alert a SOC with full telemetry context rather than one isolated
event; \textbf{isolate} the endpoint from the network to stop lateral movement while leaving an
analyst able to investigate remotely; kill the malicious process outright; quarantine a file; or
preserve forensic evidence with a memory or disk snapshot.

\begin{lstlisting}[language=Python, caption={Flagging suspicious process events by parent-child relationship and network behavior}]
SUSPICIOUS_PARENTS = {"winword.exe", "excel.exe", "outlook.exe"}
SUSPICIOUS_CHILDREN = {"powershell.exe", "cmd.exe", "wscript.exe"}

def flag_suspicious_processes(events):
    """events: list of dicts like {"pid": str, "process": str,
       "parent_process": str, "made_network_connection": bool}.
       Returns sorted PIDs showing a suspicious parent-child spawn
       combined with unexpected network activity."""
    flagged = []
    for e in events:
        spawn_suspicious = (e["parent_process"] in SUSPICIOUS_PARENTS and
                             e["process"] in SUSPICIOUS_CHILDREN)
        if spawn_suspicious and e["made_network_connection"]:
            flagged.append(e["pid"])
    return sorted(flagged)
\end{lstlisting}

\begin{misconception}
\textbf{``EDR is just better antivirus.''} Antivirus primarily checks files against known
signatures at rest. EDR monitors continuous runtime behavior across process, file, network, and
registry streams simultaneously, and can respond by isolating, killing, or quarantining -- not
merely scanning and reporting.
\end{misconception}

\section{Log Analysis}

EDR watches live behavior on one device; logs extend that visibility across time and across every
system that generates one -- operating systems, web servers, firewalls, authentication systems.

\begin{definition}{Correlation}
Linking multiple log entries together -- by the same user, the same source IP, or the same time
window -- to reveal a pattern that no single entry shows on its own.
\end{definition}

A single failed login line is entirely ordinary; anyone mistypes a password occasionally. Five
failed attempts for the same username within a ten-second window, all from the same source IP, is
a different fact entirely. No single line in that sequence looks alarming in isolation, but the
pattern itself is a textbook \textbf{brute-force} signature: rapid, clustered failures targeting one
account. Analysts formalize this into a \textbf{detection rule}: flag any account with five or
more failures within any sixty-second window. A \textbf{SIEM} platform then runs rules like this
automatically across every log source in an organization at once.

None of this works if the attacker can simply delete the evidence. An attacker with local
administrator access can clear the Windows Event Log or edit \texttt{auth.log} directly, erasing
exactly the entries that would reveal the intrusion. This is why detection programs forward logs off
the host, ideally in near-real-time, to a centralized collector or SIEM. Once a copy of an event
leaves the machine that generated it, an attacker who compromises that machine afterward can no
longer reach it. Centralized storage adds a second layer with \textbf{tamper-evident retention} --
write-once storage or cryptographically chained log entries, where altering or deleting a past entry
breaks a verifiable chain and is immediately detectable. Together, off-host forwarding and
tamper-evident retention answer a question local logging alone cannot: not just what happened, but
whether the record of what happened can still be trusted.

\begin{example}[Reading real authentication events]
On Windows, a failed logon is recorded as \textbf{Event ID 4625}, and its \texttt{Status} and
\texttt{Sub Status} fields distinguish \emph{why} it failed. Sub status \texttt{0xC000006A} means
the account name was correct but the password was wrong (a genuine guess). Sub status
\texttt{0xC0000064} means the account name itself does not exist, which is useful for spotting an
attacker enumerating usernames. \textbf{Event ID 4624} is the corresponding successful logon,
recorded on the machine the session was created on. \textbf{Event ID 4740} fires when an account is
locked out after too many failures. On a domain using Kerberos, \textbf{Event ID 4771} records a
failed Kerberos pre-authentication instead -- a Kerberos-specific failure signal distinct from 4625.
A detection rule built only around 4625 would silently miss it entirely.\srccite{4.4}{1}
\begin{center}
\ttfamily 4625 (repeated) $\rightarrow$ 4740 (lockout) $\rightarrow$ 4624 (success, new IP)
\end{center}
is the signature of a single account under sustained attack. On Linux, the equivalent evidence lives
in \texttt{/var/log/auth.log}:
\begin{lstlisting}[language=bash, caption={Linux auth.log lines showing a failed then successful SSH login}]
Sep  2 03:14:07 web01 sshd[8841]: Failed password for admin from 203.0.113.9 port 51322 ssh2
Sep  2 03:14:09 web01 sshd[8843]: Failed password for admin from 203.0.113.9 port 51330 ssh2
Sep  2 03:14:41 web01 sshd[8850]: Accepted password for admin from 203.0.113.9 port 51401 ssh2
\end{lstlisting}
No single line above is alarming by itself, but the same correlation principle applies. Repeated
failures for one account from one source, followed by success, at 3~a.m., form the pattern -- not
the final "Accepted" line read in isolation.
\end{example}

\begin{lstlisting}[language=Python, caption={Detecting a brute-force login pattern via log correlation}]
def find_bruteforce_accounts(events, threshold=5, window_seconds=60):
    """events: list of dicts like {"user": str, "timestamp": int,
       "result": "FAILED"|"SUCCESS"}.
       Returns sorted usernames with >= threshold failures inside any
       window_seconds-wide sliding window."""
    from collections import defaultdict
    by_user = defaultdict(list)
    for e in events:
        if e["result"] == "FAILED":
            by_user[e["user"]].append(e["timestamp"])

    flagged = set()
    for user, times in by_user.items():
        times.sort()
        left = 0
        for right in range(len(times)):
            while times[right] - times[left] > window_seconds:
                left += 1
            if right - left + 1 >= threshold:
                flagged.add(user)
                break
    return sorted(flagged)
\end{lstlisting}

Threshold tuning here is the same tension seen throughout this course. A threshold set too low --
flagging after two failures -- generates constant false alarms from users who simply mistyped a
password, training analysts into \textbf{alert fatigue}. A threshold set too high, on the other hand,
lets a real, deliberately slow brute-force attempt pass through undetected. Logs matter beyond
real-time detection as well: they are frequently the \emph{only} record of what an attacker did once
that attacker has already left the system. That makes them essential for \textbf{forensics} --
reconstructing an incident after the fact rather than catching it as it happens.

\begin{example}[Distinguishing a brute-force attack from a password spray]
The brute-force detector above watches one account for many failures. A \textbf{password spray}
inverts that pattern: an attacker tries one or two common passwords against \emph{many} different
accounts from the same source. The attacker deliberately stays under any single account's lockout
threshold, so the per-account brute-force rule never fires at all. Detecting a spray requires correlating by
source IP across the whole account population instead of by username:
\begin{lstlisting}[language=Python, caption={Detecting a password-spray pattern by source IP fan-out}]
from collections import defaultdict

def find_spray_sources(events, min_accounts=8, window_seconds=300):
    """events: list of dicts like {"source_ip": str, "user": str,
       "timestamp": int, "result": "FAILED"|"SUCCESS"}.
       Returns source IPs that hit >= min_accounts distinct usernames
       with a FAILED result inside any window_seconds-wide window."""
    by_ip = defaultdict(list)
    for e in events:
        if e["result"] == "FAILED":
            by_ip[e["source_ip"]].append((e["timestamp"], e["user"]))

    flagged = set()
    for ip, hits in by_ip.items():
        hits.sort()
        left = 0
        seen = defaultdict(int)
        for right in range(len(hits)):
            ts, user = hits[right]
            seen[user] += 1
            while hits[right][0] - hits[left][0] > window_seconds:
                seen[hits[left][1]] -= 1
                if seen[hits[left][1]] == 0:
                    del seen[hits[left][1]]
                left += 1
            if len(seen) >= min_accounts:
                flagged.add(ip)
                break
    return sorted(flagged)
\end{lstlisting}
This function was run against a synthetic 20-account, one-source-IP event log during preparation of
this chapter. It correctly returned the spraying IP once that IP touched eight distinct accounts
within the five-minute window, while leaving a single legitimate user's occasional mistyped password
unflagged. The distinction matters operationally: a rule tuned only to catch many failures on one
account is structurally blind to a spray, since no individual account in a spray ever crosses that
account's own threshold.
\end{example}

\begin{misconception}
\textbf{``A single failed login is proof of an attack.''} One failed login is ordinary human
error. Attacks are recognized through correlated patterns -- multiple entries linked by user, IP,
and a tight time window -- never from one log line in isolation.
\end{misconception}

\section{Indicators of Compromise}

EDR and log analysis both ultimately produce the same kind of output: individual clues that
something may be wrong.

\begin{definition}{Indicator of compromise (IOC)}
A piece of forensic evidence -- an artifact left behind on a system or network -- that suggests a
device has been breached.
\end{definition}

IOCs come from several layers. \textbf{File-based} IOCs include a file matching a known-malicious
hash or an unexpected executable in a startup directory. \textbf{Network-based} IOCs include a
connection to a known command-and-control IP or DNS lookups for algorithmically generated domain
names, a technique malware uses to find its C2 server without ever hardcoding one address.
\textbf{Host-based} IOCs include a new administrator account nobody remembers creating, a
persistence mechanism like an unowned scheduled task, or antivirus protection silently disabled.
\textbf{Behavioral} IOCs are patterns rather than single artifacts -- a login at 3~a.m.\ from an
account that has never logged in outside business hours.

\begin{example}[Why one IOC is rarely a verdict]
A new administrator account appears on a server. Alone, this could easily be an IT technician's
legitimate maintenance action, or it could be an attacker establishing persistence -- the IOC alone
cannot distinguish the two. If that same account appears \emph{alongside} a disabled antivirus
service and a connection to a known-malicious IP, the combination becomes far more convincing. Three
independent artifacts pointing toward the same conclusion form dramatically stronger evidence than
any one of them alone. That is the identical correlation principle from log analysis, applied here
to compromise evidence broadly rather than login timestamps specifically.
\end{example}

Because different IOC types carry very different severity, mature detection systems assign each a
\textbf{weight} rather than simply counting how many are present. A confirmed known-malicious file
hash is strong evidence and receives a high weight, while a handful of failed logins is comparatively
weak evidence and receives a low one. Summing every observed IOC's weight produces a
\textbf{compromise score}, compared against thresholds to reach a verdict: clean below a low
threshold, suspicious in a middle band warranting investigation, and compromised -- warranting
immediate isolation -- above a high one. That is the same weighted-scoring logic CVSS applies to
vulnerabilities in Chapter~\ref{ch:4.3}.

\begin{lstlisting}[language=Python, caption={Scoring a device's compromise level from weighted IOCs}]
IOC_WEIGHTS = {
    "known_malicious_hash": 10,
    "c2_connection": 8,
    "new_unapproved_admin": 5,
    "antivirus_disabled": 6,
    "off_hours_login": 2,
}

def compromise_verdict(observed_iocs, suspicious_threshold=6, compromised_threshold=15):
    """observed_iocs: list of IOC name strings observed on a device.
       Returns "clean", "suspicious", or "compromised"."""
    score = sum(IOC_WEIGHTS.get(ioc, 0) for ioc in observed_iocs)
    if score >= compromised_threshold:
        return "compromised"
    if score >= suspicious_threshold:
        return "suspicious"
    return "clean"
\end{lstlisting}

An \textbf{indicator of attack (IOA)} is a related but distinct concept. It focuses on an attacker's
intent and technique \emph{while} it is happening -- a process attempting to enumerate every file
with a specific extension, which is exactly what ransomware does moments before encrypting them. An
IOC, by contrast, is typically evidence left \emph{after} a step has already occurred, such as a file
with a \texttt{.locked} extension that already exists. IOAs aim to catch an
attack in progress; IOCs confirm that something already happened. Real detection programs use both,
which is the point of studying EDR, log correlation, and IOC scoring together rather than as three
unrelated tools.

\begin{misconception}
\textbf{``Any single suspicious artifact proves a device is compromised.''} One IOC alone is
usually ambiguous and often has an entirely innocent explanation. Confident verdicts come from
combining multiple independent IOCs, exactly as correlating multiple log entries does in the
previous section.
\end{misconception}

\section*{Review Questions}
\begin{reviewquestions}
\item An email macro launches PowerShell, which then opens an outbound network connection.
Explain why traditional signature-based antivirus is likely to miss this chain entirely, and
explain which EDR detection method catches it instead and why.
\item An authentication log shows five failed logins for one account within eight seconds, all
from the same IP, followed by a successful login from that same IP. Explain why reading the final
successful login in isolation is misleading, and describe the correlation that reveals the true
pattern.
\item A detection rule is tuned to flag any account with two or more failed logins in a day.
Explain the specific operational problem this creates and name the phenomenon it produces among
analysts.
\item An organization keeps its authentication logs only on each local device, with no forwarding
to a central collector. Explain the specific evidentiary weakness this creates once an attacker
gains administrative access to a device, and explain what tamper-evident retention adds even after
logs are forwarded somewhere else.
\item A server shows a new administrator account and a disabled antivirus service, but no known-
malicious network connection. Using the weighted-scoring approach from this chapter, explain why
this combination might still cross a "suspicious" threshold even without every category of IOC
present.
\item Explain the difference between an indicator of compromise and an indicator of attack, and
explain why a mature detection program needs both rather than relying on just one.
\item A security team notices that a source IP triggers no per-account brute-force alert despite
having failed logins against twenty different accounts over a few minutes. Explain why the
brute-force detector from this chapter misses this entirely, and describe the specific
correlation dimension a password-spray detector uses instead.
\item Using the EDR telemetry diagram, explain why an attacker using only living-off-the-land
techniques (legitimate tools like PowerShell, no custom malware file at all) can still be caught,
even though every individual telemetry stream -- the process list, the file system, the registry --
shows nothing that would trigger a signature match on its own.
\end{reviewquestions}

\begin{answerkey}
\item Traditional signature-based antivirus checks files against known-malicious hashes at rest;
PowerShell itself is a legitimate signed Windows binary with no malicious signature to match, so it
never triggers a file-based check. EDR's behavioral detection instead catches the chain by scoring
the unusual \emph{combination} of events -- an Office application spawning PowerShell, an unexpected
outbound network connection, a new file written to a temp directory. None of those events is damning
alone, but together they score as high-risk living-off-the-land activity.
\answerlabel{4.4.A}
\item Reading only the final successful login is misleading because, in isolation, one successful
login line looks entirely ordinary. The correlation that reveals the true pattern links the five
failed attempts to the same username, the same source IP, and a tight eight-second window. Together
those link to Event ID 4625 with sub status \texttt{0xC000006A} repeated in rapid succession, followed
by the eventual 4624 success -- a textbook brute-force signature. It only becomes visible once the
entries are linked rather than read one at a time.
\answerlabel{4.4.D}
\item A threshold of two failures in a day is set far too low: ordinary users who simply mistype a
password will routinely cross it, generating constant false alarms. The phenomenon this produces
among analysts is \textbf{alert fatigue}, where a flood of low-value alerts trains analysts to
deprioritize or ignore the detection rule entirely, which paradoxically lowers real detection
capability.
\answerlabel{4.4.C}
\item Using the weighted-scoring approach, a new unapproved admin account (weight 5) plus a
disabled antivirus service (weight 6) sums to 11. That exceeds the suspicious threshold of 6, even
though no network-based IOC (such as a C2 connection) is present. The scoring model does not require
every IOC category to be represented. It only requires the sum of observed weights to cross a
threshold -- two independent, moderately weighted artifacts together can be as convincing as one
heavily weighted one.
\answerlabel{4.4.C}
\item Local-only logs are only as trustworthy as the device that generates them: once an attacker
has administrative access, they can edit or delete the very Event ID or \texttt{auth.log} entries
that would reveal the intrusion, destroying the evidence before anyone reads it. Forwarding a copy
of each event off the host, ideally in near-real-time, removes it from the attacker's reach the
moment it is written -- the attacker would have to compromise the collector too, a categorically
harder task. Tamper-evident retention closes the remaining gap: even at the centralized store,
write-once storage or cryptographically chained entries mean that any later edit or deletion breaks
a verifiable chain and is caught, rather than trusted silently.
\answerlabel{4.4.B}
\item An indicator of compromise (IOC) is evidence that a step has already occurred, such as a file
with a \texttt{.locked} extension. An indicator of attack (IOA) focuses on an attacker's intent and
technique while it is happening, such as a process enumerating every file with a specific extension
moments before ransomware would encrypt them. A mature program needs both because IOAs aim to
interrupt an attack in progress while it can still be stopped, while IOCs confirm what already
happened and support the forensic reconstruction that follows.
\answerlabel{4.4.B}
\item The brute-force detector correlates failures by username, counting how many times one
account was attacked. A password spray deliberately stays under any single account's threshold
instead, trying only one or two passwords per account before moving to the next, so no individual
username ever accumulates enough failures to fire the rule. A password-spray detector instead
correlates by source IP across the entire account population, counting how many distinct usernames
one source touched within a time window. The fan-out across accounts is the signal, not the failure
count against any one of them.
\answerlabel{4.4.D}
\item Living-off-the-land techniques use tools that are individually legitimate. No single
telemetry stream -- the process list showing PowerShell running, the file system showing a
temp-directory write, the network stream showing an outbound connection -- contains anything a
signature check would flag, since none of those artifacts is malicious in isolation. The
behavioral model in the diagram does not evaluate any one stream alone. It scores the combination
across all of them at once -- an unusual parent-child spawn relationship, together with unexpected
network activity, together with a new file. That combined, cross-stream pattern is exactly what a
purely per-stream signature check has no way to see.
\answerlabel{4.4.B}
\end{answerkey}

\section*{Sources}
\begin{enumerate}[label=\lbrack\arabic*\rbrack, leftmargin=2.2em]
\item \srcitem{4.4}{1}Microsoft Corporation, ``4625(F) An account failed to log on'' and
``4771(F) Kerberos pre-authentication failed,'' \emph{Windows Security Auditing Events} reference,
Microsoft Learn -- documenting the failure sub status codes \texttt{0xC000006A} (wrong password) and
\texttt{0xC0000064} (account does not exist), and confirming Event ID 4771 is logged only on domain
controllers.
\end{enumerate}
